% Correspondencia editorial: gestion/mapa_transportes_auxiliares.md.
% Realizaciones auxiliares: no sustituyen el eje de compatibilidad y restitución.
\section{Transport compensé, mémoire géométrique et confrontation énergétique}
\label{md:aux:transporte}

\subsection{Origine opératorielle et domaine de la réalisation}
\label{md:aux:origen}

La construction conserve APP--TRIT--TPK et l’état enrichi. APP
réunit les feuillets arithmétiques avec leurs résidus et quotients ; TRIT conserve
le régime et l’orientation ; TPK compose sélection, transport et
mise à jour de la mémoire. Les représentations matricielles utilisées
ci-dessous sont des réalisations ultérieures. L’holonomie nonadique
$\Gamma_9$ et sa mémoire maintiennent leurs propres opérateurs.

Les antécédents employés sont les représentations quadratiques
elliptique, hyperbolique et parabolique du TRIT, la différentielle des
transports inversibles, la fonctionnelle énergétique
$E=\langle Df,WDf\rangle$ avec $W>0$, le retour avec mémoire
translationnelle, l’ellipse d’action avec produit des demi-axes $2\hbar$
et sa réalisation LC, la polarisation des vacances et le couplage
réversible de deux modes. Le retour translationnel et le commutateur
qui suit sont des objets distincts. Le lacet elliptique--parabolique et
son bilan d’aire orientée seront ensuite utilisés comme collection
de confrontation.

La composition compensée et le protocole de confrontation réunissent ces
réalisations. La théorie algébrique des traces de commutateurs dans
$\operatorname{SL}(2)$ constitue un antécédent mathématique
établi.\footnote{Goldman, \emph{Trace Coordinates on
Fricke spaces}, \href{https://arxiv.org/abs/0901.1404}{arXiv:0901.1404}.}
L’objet de cet appendice est leur composition concrète avec les
lecteurs déclarés. Les paramètres $a,b$ sont des paramètres des
réalisations, non des constantes physiques ajustées.

La portée s’arrête à cette réalisation locale. L’identification de
chemins TPK admettant un dispositif matériel avec ces paramètres
exige de composer leurs lecteurs effectifs. Une matrice inversible
arbitraire n’atteste pas à elle seule un chemin TPK, et l’identification
du défaut de commutation à la torsion de Cartan nécessite sa
connexion et sa soudure.

\subsection{Compensation des opérations et conservation de l’aire}
\label{md:aux:conmutador}

Les matrices de régime sont
\begin{equation}
 J_{\mathrm e}=\begin{pmatrix}0&1\\-1&0\end{pmatrix},\qquad
 J_{\mathrm h}=\begin{pmatrix}0&1\\1&0\end{pmatrix}.
 \label{md:aux:regimenes}
\end{equation}
Elles satisfont $J_{\mathrm e}^2=-I$, $J_{\mathrm h}^2=I$ et
$[J_{\mathrm e},J_{\mathrm h}]=2\operatorname{diag}(1,-1)$.
On définit
\begin{equation}
 R(a)=e^{aJ_{\mathrm e}},\qquad B(b)=e^{bJ_{\mathrm h}},\qquad
 C_{\mathrm{aux}}(a,b)=R(a)B(b)R(a)^{-1}B(b)^{-1}.
 \label{md:aux:def-compensado}
\end{equation}
L’action sur les colonnes se lit de droite à gauche. Le protocole
contient chaque transformation et son inverse, mais la suite n’est pas
le retour en arrière complet du parcours. Un retour en arrière complet aurait,
par exemple, la forme $RBB^{-1}R^{-1}=I$. Le symbole
$C_{\mathrm{aux}}$ distingue ce commutateur elliptique--hyperbolique
du contre-angle $C^*$ et du lacet de formes conjuguées
$H_{\mathrm{lazo}}$ étudié dans le corps principal.

\begin{proposition}
\label{md:aux:prop-traza}
Pour $a,b\in\mathbb R$,
\begin{equation}
 \det C_{\mathrm{aux}}=1,\qquad
 \operatorname{tr}C_{\mathrm{aux}}=2+4\sin^2a\sinh^2b.
 \label{md:aux:traza-compensado}
\end{equation}
\end{proposition}
\begin{proof}
Posons $c_a=\cos a$, $s_a=\sin a$, $u_b=\cosh b$ et $v_b=\sinh b$.
Alors
\[
 R=\begin{pmatrix}c_a&s_a\\-s_a&c_a\end{pmatrix},\qquad
 B=\begin{pmatrix}u_b&v_b\\v_b&u_b\end{pmatrix},
\]
avec $c_a^2+s_a^2=1$ et $u_b^2-v_b^2=1$. Toutes les matrices ont un
déterminant égal à un. La multiplication donne
\begin{equation}
 \begin{aligned}
 (C_{\mathrm{aux}})_{11}&=u_b^2+2c_as_au_bv_b-(c_a^2-s_a^2)v_b^2,\\
 (C_{\mathrm{aux}})_{12}&=-2s_a^2u_bv_b-2c_as_av_b^2,\\
 (C_{\mathrm{aux}})_{21}&=-2s_a^2u_bv_b+2c_as_av_b^2,\\
 (C_{\mathrm{aux}})_{22}&=u_b^2-2c_as_au_bv_b-(c_a^2-s_a^2)v_b^2.
 \end{aligned}
 \label{md:aux:entradas-compensado}
\end{equation}
La somme diagonale est
$2u_b^2-2(c_a^2-s_a^2)v_b^2=2+4s_a^2v_b^2$, comme affirmé.
\end{proof}

Lorsque $s_av_b\ne0$, la trace dépasse deux et le commutateur est
hyperbolique, avec des valeurs propres positives réciproques
\begin{equation}
 e^{\lambda},\ e^{-\lambda},\qquad
 \lambda=2\operatorname{arsinh}|\sin a\sinh b|>0.
 \label{md:aux:espectro-compensado}
\end{equation}
Cela s’obtient à partir de $\cosh\lambda=\operatorname{tr}(C_{\mathrm{aux}})/2$
et $\cosh(2\operatorname{arsinh}x)=1+2x^2$.
Si $s_av_b=0$, alors $b=0$ ou $R=\pm I$, et
$C_{\mathrm{aux}}=I$.
Le développement local est
\begin{equation}
 C_{\mathrm{aux}}=I+2ab\operatorname{diag}(1,-1)
 +O(|a|^2|b|+|a||b|^2).
 \label{md:aux:expansion-compensado}
\end{equation}
Le résidu commence par le produit des deux opérations. Appliquer
l’une quelconque d’entre elles suivie de sa propre inverse produit l’identité.

Une matrice réelle bidimensionnelle de déterminant un conserve la forme
d’aire. Ainsi, $C_{\mathrm{aux}}$ transforme l’ellipse
d’action en une autre ellipse d’aire $h$. Si $\Sigma>0$ est une matrice
de covariance,
\begin{equation}
 \Sigma'=C_{\mathrm{aux}}\Sigma C_{\mathrm{aux}}^{\mathsf T},
 \qquad \det\Sigma'=\det\Sigma.
 \label{md:aux:covarianza-area}
\end{equation}
La distribution des demi-axes et leur orientation peuvent changer.
Les valeurs propres $e^{\pm\lambda}$ décrivent des directions propres et des
itérations ; les valeurs singulières déterminent une autre lecture et, en
général, ne coïncident pas avec elles. La variation des demi-axes euclidiens
en une seule application nécessite la métrique correspondante.

La conservation de l’aire d’action est compatible avec une
accumulation de déformation orientée. L’observable géométrique
$\lambda$ dépend de $a,b$ et ne constitue pas une constante universelle
supplémentaire. La conservation du déterminant ou de l’aire, à elle
seule, n’établit pas la conservation de l’énergie physique.

\subsection{Lecteur positif du résidu et inversion du parcours}
\label{md:aux:residuo}

Dans une fibre de retour, avec un champ $f$ et un poids hermitien $W>0$
déclarés, le lecteur énergétique prend la forme
\begin{equation}
 E_C(f)=\langle(I-C_{\mathrm{aux}})f,
                  W(I-C_{\mathrm{aux}})f\rangle.
 \label{md:aux:energia-residuo}
\end{equation}
Pour $s_av_b\ne0$, le commutateur n’a pas de valeur propre égale à un et
$E_C(f)>0$ pour tout $f\ne0$. Le retour en arrière complet produit un résidu
nul. La normalisation dimensionnelle appartient à la fonctionnelle de
départ : ses coordonnées numériques nécessitent leur réalisation avant
d’être interprétées comme des joules, de la chaleur, une masse ou un courant de conduction.

L’inversion complète du parcours inverse $C_{\mathrm{aux}}$.
La trace et $\lambda$ identifient l’opérateur et son inverse. La
restitution de direction nécessite une lecture supplémentaire : le vecteur
résiduel, les composantes du transport ou une sonde avec résolution
de phase. Cette comparaison limite concrètement la suffisance
d’une lecture exclusivement scalaire.

Un exemple rationnel, indépendant de données physiques, est
\begin{equation}
 R=\begin{pmatrix}3/5&4/5\\-4/5&3/5\end{pmatrix},\qquad
 B=\begin{pmatrix}5/4&3/4\\3/4&5/4\end{pmatrix}.
 \label{md:aux:ejemplo-racional}
\end{equation}
On y obtient $\operatorname{tr}C_{\mathrm{aux}}=86/25$,
$\det C_{\mathrm{aux}}=1$ et les valeurs propres
$(43\pm6\sqrt{34})/25$. Les contrôles rationnels du supplément
comprennent cet exemple et $169$ paires, avec des cas nuls, un retour en arrière
complet et la conservation du déterminant de covariance. Ce sont
des vérifications finies des identités démontrées.

\subsection{Couplage réversible et lecture réduite de l’intrication}
\label{md:aux:acoplamiento}

Pour $H\in\operatorname{Sp}(2,\mathbb R)$, la réalisation à deux
modes utilise le couplage
\begin{equation}
 U_H=\frac19\begin{pmatrix}
 8I+H&\sqrt8(H-I)\\
 \sqrt8(H-I)&I+8H
 \end{pmatrix}.
 \label{md:aux:acoplamiento-uh}
\end{equation}
Avec deux gaussiennes pures identiques préparées à partir de l’ellipse
d’action et avec un transport conjoint de la mise en coordonnées initiale, le lecteur
réduit satisfait
\begin{equation}
 \nu(H)^2=1+\frac8{81}
    \bigl\{\operatorname{tr}(HH^{\mathsf T})-2\bigr\},
 \qquad\operatorname{Tr}\rho_{\mathrm{vis}}^2=\nu(H)^{-1}.
 \label{md:aux:mezcla-generica}
\end{equation}
La préparation et les lecteurs restent fixés. La
transformation conjointe est réversible et conserve la pureté globale.

\begin{corollary}
\label{md:aux:cor-mezcla}
Pour $C_{\mathrm{aux}}(a,b)$,
\begin{equation}
 \operatorname{tr}(C_{\mathrm{aux}}C_{\mathrm{aux}}^{\mathsf T})-2
 =16\sin^2a\sinh^2b\cosh^2b,
 \label{md:aux:traza-cuadratica}
\end{equation}
et
\begin{equation}
 \boxed{\nu_C^2=1+\frac{32}{81}\sin^2a\sinh^2(2b).}
 \label{md:aux:mezcla-compensado}
\end{equation}
\end{corollary}
\begin{proof}
Les quatre entrées de \eqref{md:aux:entradas-compensado}
s’écrivent $d+k$, $l-m$, $l+m$, $d-k$, où
\[
 d=1+2s_a^2v_b^2,\quad k=2c_as_au_bv_b,\quad
 l=-2s_a^2u_bv_b,\quad m=2c_as_av_b^2.
\]
La somme de leurs carrés est
$2(d^2+k^2+l^2+m^2)=2+16s_a^2u_b^2v_b^2$.
La substitution dans \eqref{md:aux:mezcla-generica} et
$\sinh(2b)=2\sinh b\cosh b$ prouve le résultat.
\end{proof}
L’exemple \eqref{md:aux:ejemplo-racional} produit
$\nu_C^2=17/9$ et la pureté $3/\sqrt{17}$.

L’annulation de $\hbar$ et de la déformation initiale utilise
la préparation déclarée et le transport conjoint de la mise en coordonnées.
Un protocole avec $s_av_b\ne0$ produit un mélange réduit ;
l’inversion complète du couplage peut le défaire. Mélange
réduit, énergie stockée et chaleur conservent leurs fonctions
comme lectures distinctes.

\subsubsection{Contraste signé dans une collection elliptique--parabolique}
\label{md:aux:familia-hn}

La collection de lacets construite est
\begin{equation}
 H_n=\begin{pmatrix}1+n&n^2\\n&n^2-n+1\end{pmatrix},
 \qquad n\in\mathbb Z.
 \label{md:aux:hn}
\end{equation}
Son bilan alterné de mémoire dépend de $n^2$, tandis que le
quotient des déterminants de covariance est
\begin{equation}
 \mathcal R_n=\frac{\det V_{\mathrm{vis},n}}{\det V_0}
 =1+\frac8{81}n^2(2n^2-2n+5).
 \label{md:aux:covarianza-hn}
\end{equation}
En soustrayant les deux orientations, les termes pairs s’annulent :
\begin{equation}
 \boxed{\mathcal R_{-n}-\mathcal R_n=\frac{32n^3}{81}.}
 \label{md:aux:diferencia-direccional}
\end{equation}
Pour $n=1$ et $n=-1$, les réponses sont $121/81$ et $153/81$,
avec les puretés $9/11$ et $3/\sqrt{17}$, respectivement.
Cette collection satisfait $H_{-n}\ne H_n^{-1}$ en général ;
elle n’est pas non plus l’itération $H_1^n$. Le changement du paramètre
orienté doit être distingué du retour en arrière complet des
opérations. Sa lecture alternée d’aire conserve une
identification distincte de celle de la fonctionnelle de Barbero.

Le critère de mélange emploie le commutateur $[D,J]$, avec
$D=\sqrt8(H-I)/9$. Si $D$ commute avec la structure complexe
compatible avec l’état initial, il peut transporter de la mémoire
sans mélanger ce produit gaussien. S’il ne commute pas, le lecteur
détermine le mélange. L’absence de mélange pour cet état
est une propriété distincte de l’absence de dissipation matérielle.

La différence \eqref{md:aux:diferencia-direccional} est un corollaire
de la collection et de son lecteur, non une constante universelle.
La formule de $C_{\mathrm{aux}}$ compose le même lecteur avec
une autre collection admise par le théorème. Une confrontation matérielle
doit déterminer les transports et les couplages avant de mesurer ;
la programmation arbitraire de ces matrices atteste leur
émulation et conserve une portée différente.

\subsection{Deuxième moment du courant de vacances}
\label{md:aux:vacancias}

Le lecteur de capacité fixe
\[
 \delta=\log_{10}(10/9),\qquad z_n\in\{0,1\},\qquad
 P_N=\sum_{n<N}(z_n-\delta),\qquad |P_b-P_a|<1.
\]
Sur des intervalles égaux $t_0$, le courant centré est
$I_n=A_I(z_n-\delta)$, avec $A_I=u_Q/t_0$.
Le symbole $A_I$ distingue cette amplitude de l’angle principal.
Pour $L=b-a$,
\begin{equation}
 \frac1{A_I^2}\sum_{n=a}^{b-1}I_n^2
 =L\delta(1-\delta)+(1-2\delta)(P_b-P_a).
 \label{md:aux:segundo-momento}
\end{equation}
La preuve consiste à développer $(z_n-\delta)^2$, à utiliser
$z_n^2=z_n$ et à sommer. Le déséquilibre signé demeure
borné, tandis que
$I_{\mathrm{rms}}/|A_I|\longrightarrow\sqrt{\delta(1-\delta)}$.
L’amplitude $A_I$ est le saut entre les deux niveaux de courant.
La confrontation conserve la réalisation temporelle et doit comptabiliser
tout couplage qui filtre la suite.

Un bilan signé nul peut coexister avec une lecture
quadratique positive. Si l’on ajoute une résistance idéale $R_E>0$
et que l’on maintient ce courant par intervalles, le modèle ultérieur
de confrontation donne
\begin{equation}
 Q_{R_E}=R_Et_0\sum_n I_n^2.
 \label{md:aux:contraste-resistivo}
\end{equation}
La résistance appartient à cette réalisation ultérieure. Sa
fonction consiste à confronter une réduction possible de chaleur
à un bilan explicite.

\subsection{Conditions d’une confrontation matérielle}
\label{md:aux:contraste-material}

Le protocole sépare deux questions. Pour la mémoire d’ordre,
on compare des parcours compensés, des échanges d’ordre et des
retours en arrière complets au moyen d’une réponse rémanente vectorielle
ou résolue en phase. Une intensité ou une charge nette isolées
peuvent perdre la différence recherchée. Pour la récupération
d’énergie, on mesure l’entrée, l’énergie utile ou restituée, la variation
du stockage et la chaleur, y compris la préparation et la restauration
du dispositif. Le retour de phase conserve séparée la
variation de l’énergie stockée.

Dans une confrontation scalaire à un port, la relation
$v_2(t)=v_1(T-t)$ conserve la distribution des amplitudes et le
spectre de puissance. En régime périodique, un même système
linéaire stationnaire prédit une puissance moyenne identique. Une loi
instantanée $i=g(v)$ conserve également l’intégrale de $vi$.
Une différence reproductible examinerait une dépendance historique,
mais son attribution nécessite de comparer l’hystérésis, l’échauffement
et des modèles dynamiques alternatifs. Avec plusieurs ports, il faut
en outre conserver les spectres croisés.

Une réduction de la dissipation est testée à transfert
utile, fidélité et rythme égaux. L’atténuation doit distinguer l’énergie,
l’amplitude du signal, la population de porteurs et la charge conservée
avant de lui attribuer une origine gravitationnelle. La distinction entre
stockage, récupération et pertes possède des précédents
expérimentaux dans la récupération électrocalorique.\footnote{Defay
et al., \emph{Enhanced electrocaloric efficiency via energy
recovery} (2018),
\href{https://doi.org/10.1038/s41467-018-04027-9}{doi:10.1038/s41467-018-04027-9}.
Cet antécédent ne constitue pas une preuve expérimentale du
transport HMT composé ici.}

Le résultat de l’appendice est une identité générale au sein
de la collection matricielle spécifiée, un lecteur positif du
résidu, un contrôle d’inversion et une loi exacte du deuxième
moment du courant. Sa portée distingue ces résultats
d’une nouvelle constante universelle, d’un prototype sans échauffement
ou d’une démonstration de désintégration électronique gravitationnelle.
Le lien prospectif avec un dispositif doit être fixé
avant d’obtenir les données de confrontation.

\section{Action, mémoire réduite et travail récupérable}
\label{md:aux:trabajo}

\subsection{Composition de l’action avec une réalisation énergétique}
\label{md:aux:accion}

La composition reçoit les sorties d’action et les transports
de mémoire d’APP--TRIT--TPK, en conservant l’état enrichi
comme antécédent. Son domaine est la réalisation gaussienne et
énergétique déclarée. La passivité et l’ergotropie sont des concepts
établis ; ils sont employés pour étudier le couplage spécifique
et les lecteurs déjà construits.

L’équation d’action maintient ses coefficients explicites :
\begin{equation}
 H_5=\frac12\sqrt{\frac{1000\alpha}{\varphi}}-\alpha
 +\frac9{16}\alpha^2-\frac59\alpha^3
 +\frac7{48}\alpha^4-\frac1{54}\alpha^5,
 \label{md:aux:h5}
\end{equation}
\begin{equation}
 D_A=e^{-100\pi\alpha/9},\qquad
 \mathcal C_\pi=169\alpha^6+\frac{D_A\alpha^7}{1-D_A\alpha},
 \label{md:aux:correccion-accion}
\end{equation}
\begin{equation}
 \begin{aligned}
 \hbar_{\mathrm{pre}}&=s_0H_5,\\
 \hbar_{\mathrm{ret}}&=s_0\left(H_5-\frac{90}{\pi}\mathcal C_\pi\right),
 \qquad h_s=2\pi\hbar_s,\qquad s_0=10^{-34}\mathcal U_S.
 \end{aligned}
 \label{md:aux:secciones-accion}
\end{equation}
La décennie et la mise en coordonnées dimensionnelle conservent des antécédents
distincts. La formule maintient les coordonnées $\pi,\varphi$,
l’orientation et l’échelle du lecteur, en plus de $\alpha$.
L’égalité dimensionnelle conserve son unité d’action.

Une section positive étant fixée, on écrit $\hbar=\hbar_s$.
L’ellipse a pour demi-axes $\sqrt{2\hbar}e^{\zeta/2}$ et
$\sqrt{2\hbar}e^{-\zeta/2}$, où
$\zeta=\operatorname{artanh}(C^*/A)$. Son aire orientée satisfait
\begin{equation}
 \mathcal A(\theta)=\hbar\theta,\qquad
 \mathcal A(2\pi)=h.
 \label{md:aux:area-fase}
\end{equation}
Dans la réalisation LC, avec une horloge $\omega$ et une impédance de
référence déclarées,
\begin{equation}
 \mathcal H_\zeta(Q,P)
 =\frac\omega2\left(e^{-\zeta}Q^2+e^\zeta P^2\right).
 \label{md:aux:hamiltoniano-lc}
\end{equation}
Sur l’orbite, on obtient $\mathcal H_\zeta=\hbar\omega=hf$,
avec $f=\omega/(2\pi)$. La convention hamiltonienne donne
$\dot\theta=-\omega$ et, par conséquent,
$d\mathcal A/dt=-\hbar\omega$. L’énergie positive coïncide
avec la grandeur de cette vitesse aréolaire dans la réalisation et
la section fixe considérées. L’aire accumulée et l’action
lagrangienne d’une trajectoire arbitraire sont des objets distincts.

L’équation interne détermine l’échelle d’action par unité
de phase ; l’horloge fixe son rythme. La déformation de l’ellipse
et l’ordre du transport conservent une information supplémentaire par rapport à
$hf$. La préparation gaussienne employée ensuite est distincte
d’un point de l’orbite classique : son énergie initiale est
$hf/2$, au lieu de $hf$.

\subsection{Préparation et lecteurs du résultat énergétique}
\label{md:aux:hipotesis-gaussianas}

On prépare deux modes gaussiens purs, centrés et identiques,
avec la covariance
\begin{equation}
 V_0=\frac\hbar2MM^{\mathsf T},\qquad
 M=\operatorname{diag}(e^{\zeta/2},e^{-\zeta/2}).
 \label{md:aux:preparacion}
\end{equation}
$U_H$ de \eqref{md:aux:acoplamiento-uh} agit, transporté vers
la mise en coordonnées initiale au moyen de $M\oplus M$. L’état conjoint
demeure pur. Les lecteurs énergétiques finaux ont une
même fréquence $\omega$, conservent la métrique initiale
$G_\zeta=M^{-\mathsf T}M^{-1}$ et ne comprennent pas d’énergie
d’interaction résiduelle entre les modes. Les opérations
d’extraction commencent et se terminent avec ces mêmes lecteurs.
L’accessibilité physique des opérations admises est une
hypothèse de la borne de travail, distincte de la construction
d’un dispositif qui les réalise.

On définit
\begin{equation}
 \begin{gathered}
 \tau=\operatorname{tr}(HH^{\mathsf T})-2\geq0,\\
 W_v=\frac{8I+HH^{\mathsf T}}9,\qquad
 W_m=\frac{I+8HH^{\mathsf T}}9.
 \end{gathered}
 \label{md:aux:covarianzas-reducidas}
\end{equation}
Les covariances sont $V_j=(\hbar/2)MW_jM^{\mathsf T}$.
La pureté visible est notée $\mu=\operatorname{Tr}\rho_v^2$
et sa valeur propre symplectique normalisée $\nu=\mu^{-1}$.
Dans cet appendice, $\nu$ remplit cette fonction ; la fréquence
ordinaire s’écrit $f$ et la fréquence angulaire $\omega$.

\subsection{Identité entre énergie visible et pureté}
\label{md:aux:energia-pureza}

\begin{proposition}
\label{md:aux:prop-energia-pureza}
Avec la préparation et les lecteurs précédents, l’énergie
visible au-dessus de l’état initial satisfait
\begin{equation}
 \frac{\Delta E_v}{hf}=\frac\tau{36},\qquad
 \nu^2=1+\frac{8\tau}{81},\qquad
 \boxed{\mu^{-2}-1=\frac{32}{9}\frac{\Delta E_v}{hf}.}
 \label{md:aux:identidad-energia-pureza}
\end{equation}
\end{proposition}
\begin{proof}
L’état est centré et le lecteur quadratique demeure fixe.
Par conséquent,
\[
 E_v=\frac\omega2\operatorname{tr}(G_\zeta V_v)
 =\frac{\hbar\omega}{4}\operatorname{tr}W_v.
\]
Comme $\operatorname{tr}W_v=2+\tau/9$, en soustrayant
$\hbar\omega/2$, on obtient la première égalité.
La covariance réduite satisfait
$\det W_v=1+8\tau/81$ et $\mu=(\det W_v)^{-1/2}$.
L’élimination de $\tau$ prouve l’identité encadrée.
La démonstration ne nécessite ni la symétrie de $H$ ni la coïncidence
entre ses valeurs propres et ses valeurs singulières.
\end{proof}

La loi élimine les paramètres particuliers du parcours.
Le coefficient $32/9$ provient de la répartition du couplage
et l’échelle $h$ est conservée jusqu’à la normalisation.
Pour vérifier la spécificité de cette répartition, considérons
$W_p=pI+(1-p)HH^{\mathsf T}$, avec $0<p<1$. On obtient
\begin{equation}
 \mu^{-2}-1=4p\frac{\Delta E}{hf}.
 \label{md:aux:reparto-general}
\end{equation}
Le choix déclaré est $p=8/9$. Lorsque $\Delta E>0$,
une confrontation peut estimer
\begin{equation}
 p_{\mathrm{obs}}=
 \frac{\mu^{-2}-1}{4\Delta E/(hf)}
 \label{md:aux:reparto-observado}
\end{equation}
et le comparer à $8/9$ sans ajuster $p$ à ces mêmes données.
Programmer un mélangeur avec cette répartition vérifie une
émulation ; l’attribution de la répartition à un mécanisme naturel
conserve sa preuve indépendante.

L’identité utilise des gaussiennes centrées et la préparation
spécifiée. Un déplacement cohérent ajoute de l’énergie sans
changer la pureté. Le bruit thermique initial et les changements de
métrique modifient la relation. Ces variations délimitent
l’énoncé et fournissent des critères de réfutation d’une extension à des états
arbitraires.

La combinaison convexe $W_v$ est une combinaison de matrices
de covariance produite en réduisant un état gaussien conjoint.
Un mélange probabiliste de deux états purs constitue
un autre objet, généralement non gaussien ; sa pureté n’est pas
donnée par le déterminant de cette covariance.

\subsection{Travail récupérable sous opérations locales}
\label{md:aux:ergotropia}

L’ergotropie est la diminution maximale d’énergie par des
opérations unitaires cycliques avec un lecteur énergétique fixe.
La définition ne comprend pas automatiquement le coût net de
construction, de préparation ou de contrôle de l’appareil.

Le spectre réduit possède les populations
\begin{equation}
 \lambda_j=(1-r)r^j,\qquad r=\frac{\nu-1}{\nu+1}.
 \label{md:aux:espectro-gaussiano}
\end{equation}
Elles sont ordonnées de manière décroissante avec l’énergie de l’oscillateur.
L’état passif a ainsi l’énergie $E_{v,\mathrm{pas}}=hf\nu/2$.
Une opération gaussienne l’atteint :
\begin{equation}
 S=\sqrt\nu\,W_v^{-1/2},\qquad
 \det S=1,\qquad SW_vS^{\mathsf T}=\nu I.
 \label{md:aux:pasivacion}
\end{equation}
L’ordre des populations démontre également la minimalité
parmi toutes les opérations unitaires, au-delà des gaussiennes.

Par conséquent,
\begin{equation}
 \boxed{
 \frac{\mathcal W_v}{hf}
 =\frac{(\nu-1)(9\nu-7)}{32},\qquad
 \frac{E_{v,\mathrm{pas}}-hf/2}{hf}=\frac{\nu-1}{2}.}
 \label{md:aux:trabajo-visible}
\end{equation}
Les deux quantités ont pour somme
$\Delta E_v/(hf)=9(\nu^2-1)/32$. Le terme passif est une
énergie inaccessible sous les opérations locales déclarées ;
sa présence n’équivaut pas à un flux de chaleur réalisé.

\subsubsection{Accès conjoint aux corrélations}
\label{md:aux:acceso-conjunto}

Le deuxième mode satisfait $\det W_m=\det W_v=\nu^2$ et
$\Delta E_m=8\Delta E_v$. Par conséquent,
\begin{equation}
 \Delta E_{\mathrm{tot}}=9\Delta E_v=hf\frac\tau4.
 \label{md:aux:energia-conjunta}
\end{equation}
L’inversion conjointe du couplage restitue le produit
initial. L’état total est pur et ce produit est l’état
fondamental du lecteur conjoint ; l’accès conjoint permet de
récupérer idéalement $\Delta E_{\mathrm{tot}}$.
Les opérations locales indépendantes laissent deux états
passifs, chacun d’énergie $hf\nu/2$. On en déduit
\begin{equation}
 \boxed{\mathcal W_{\mathrm{conjunta}}-\mathcal W_{\mathrm{local}}
 =hf(\nu-1)=hf(\mu^{-1}-1).}
 \label{md:aux:ventaja-conjunta}
\end{equation}

L’identité quantifie le travail supplémentaire accessible
au moyen des corrélations. La préparation initiale a nécessité
de l’énergie et le bilan d’un cycle complet comprend cet apport.
La relation compare des lectures énergétiques et des opérations sur
l’état, sans identifier l’information et l’énergie à une
même grandeur physique.

Dans cet énoncé, ``local'' signifie des opérations unitaires de produit
$U_v\otimes U_m$ ; les mesures,
la communication classique et la rétroaction sont exclues de cette classe. L’accroissement
provient de la restriction opérationnelle sur l’état
corrélé et de ses spectres marginaux, non d’une
énergie d’interaction résiduelle. La répartition $1:8$ concerne
les énergies au-dessus du vide, non les énergies
totales ni les ergotropies.

\subsection{Cas exacts et orientation du transport}
\label{md:aux:casos-energia}

Pour $H_n$ de \eqref{md:aux:hn},
\begin{equation}
 \tau_n=n^2(2n^2-2n+5).
 \label{md:aux:tau-hn}
\end{equation}
Dans $n=1$, on obtient
\begin{equation}
 \begin{gathered}
 \nu=11/9,\qquad\mu=9/11,\qquad
 \Delta E_v=\frac5{36}hf,\\
 \mathcal W_v=\frac1{36}hf,\qquad
 E_{v,\mathrm{pas}}-hf/2=\frac19hf.
 \end{gathered}
 \label{md:aux:caso-positivo}
\end{equation}
L’ensemble possède $\Delta E_{\mathrm{tot}}=5hf/4$, le travail
local $37hf/36$ et l’avantage conjoint $2hf/9$.
Le cas $n=-1$ partage les valeurs propres de $H_1$,
mais possède $\tau=9$ et
\begin{equation}
 \nu=\frac{\sqrt{17}}3,\qquad
 \Delta E_v=\frac{hf}{4},\qquad
 \mathcal W_v=\frac{9-2\sqrt{17}}{12}hf.
 \label{md:aux:caso-negativo}
\end{equation}
L’énergie visible ajoutée diffère dans le rapport $9/5$ à
fréquence et section d’action égales. Le changement $n\mapsto-n$
n’est pas le retour en arrière complet : $H_{-n}\ne H_n^{-1}$ en général.
La véritable inversion matricielle conserve $\tau$ en
dimension deux. L’asymétrie paramétrique de cette collection
conserve ainsi sa distinction par rapport à un renversement complet.

Pour $C_{\mathrm{aux}}$,
$\tau_C=4\sin^2a\sinh^2(2b)$. La composition énergétique donne
\begin{equation}
 \boxed{\Delta E_v(C_{\mathrm{aux}})
 =\frac{hf}{9}\sin^2a\sinh^2(2b),}
 \label{md:aux:energia-compensado}
\end{equation}
\begin{equation}
 \mathcal W_{\mathrm{conjunta}}-\mathcal W_{\mathrm{local}}
 =hf\left(\sqrt{1+\frac{32}{81}\sin^2a\sinh^2(2b)}-1\right).
 \label{md:aux:trabajo-compensado}
\end{equation}
La connexion avec Planck demeure explicite : dans la section
précédente, $hf=2\pi s_0H_5(\alpha,\varphi)f$ ; dans celle du retour,
$H_5$ est remplacé par $\eta_{\mathrm{ret}}$. La composition
maintient fixée la section utilisée ; elle ne suppose pas une
variation expérimentale de $\alpha$ ni une alternance libre
entre sections dimensionnelles.

\subsection{Boltzmann et température équivalente de l’état passif}
\label{md:aux:temperatura}

Pour $\nu>1$, l’état passif a pour température équivalente
\begin{equation}
 k_BT_{\mathrm{pas}}
 =\frac{hf}{\log[(\nu+1)/(\nu-1)]}.
 \label{md:aux:temperatura-pasiva}
\end{equation}
Lorsque l’on utilise l’horloge du transducteur,
$hf=k_B\Theta_{\mathrm{clk}}\log3$, il en résulte
\begin{equation}
 \frac{T_{\mathrm{pas}}}{\Theta_{\mathrm{clk}}}
 =\frac{\log3}{\log[(\nu+1)/(\nu-1)]}.
 \label{md:aux:razon-temperaturas}
\end{equation}
Pour $n=1$, on retrouve $\log3/\log10$. L’entropie est
celle du spectre gaussien, avec l’occupation $(\nu-1)/2$.
La composition énergétique conserve cette lecture thermique
préalablement construite. La température équivalente de
l’état passif n’implique pas que l’état réduit antérieur
ait été à l’équilibre ni que de la chaleur ait été transférée.

\subsection{Spécificité et conditions de confrontation}
\label{md:aux:contraste-trabajo}

La relation générale entre cohérence, corrélations et
extraction de travail appartient à la littérature existante.
La composition examinée spécifie une loi conjointe :
coefficient $32/9$, répartition énergétique $1:8$, collection
orientée $H_n$, échelle d’action et différence entre
accès local et conjoint.

On distingue deux niveaux expérimentaux. L’émulation
réalise $U_H$ et vérifie ses identités ; elle vérifie cette
réalisation du modèle. La prédiction physique exige de
déterminer prospectivement, à partir de la construction, le
dispositif, les degrés de liberté, la fréquence, la
préparation et le couplage qui réalisent ses objets,
sans ajuster $8/9$ aux données de la confrontation. On mesure
séparément l’énergie, la covariance, la pureté, les coûts de
contrôle et l’énergie de restauration.

Une confrontation indépendante obtient la température et la pureté
par des lecteurs indépendants de la loi qu’elle entend
vérifier. Le travail récupéré est comparé à un
bilan comprenant la préparation, les contrôles et la restauration
de la mémoire.

Les comparaisons externes comprennent des études
d’ergotropie gaussienne et de passivation,\footnote{Kua, Serafini
et Genoni, \emph{Daemonic ergotropy of Gaussian quantum
states and the role of measurement-induced purification
via general-dyne detection},
\href{https://arxiv.org/abs/2506.22288}{arXiv:2506.22288}.
Leur rôle est la comparaison avec des formules générales
d’énergie, de pureté et de passivité ; elles ne valident ni le couplage HMT
ni ses constantes.}
l’extraction d’ergotropie cohérente avec des coûts
énergétiques explicites,\footnote{\emph{Experimental
Extraction of Coherent Ergotropy and Its Energetic Cost
in a Superconducting Qubit},
\href{https://arxiv.org/abs/2506.16881}{arXiv:2506.16881}.}
et le contrôle cohérent entre les cycles de machines
quantiques.\footnote{Hou et al., \emph{Combining energy
efficiency and quantum advantage in cyclic machines},
\href{https://doi.org/10.1038/s41467-025-60179-5}{doi:10.1038/s41467-025-60179-5}.
L’expérience citée ne mesure pas le transport nonadique
composé dans cet appendice.}

Le résultat est une composition algébrique avec une
structure de départ explicite. Ses preuves fixent des
identités et des restrictions opérationnelles. Un avantage
expérimental observé, une nouvelle constante universelle,
une violation de la thermodynamique ou un mécanisme matériel
de transport sans chaleur nécessiteraient leurs démonstrations
et confrontations respectives ; ce ne sont pas des conclusions des
identités auxiliaires établies ici.
